\clearpage
\VerdierPhysicalPageAnchor{0070}%
\VerdierSemanticAnchor{ch1-2-2-3-continuation}%

\VerdierSemanticAnchor{ch1-2-2-3-eq3}
\begin{equation}
  \begin{aligned}
    \left(\int_{\varphi}^{\Pi}Y\right)^{\tau}
      &=\prod_{\varphi(\sigma)=\tau}Y^{\sigma},\\
    d_{\int_{\varphi}^{\Pi}Y}^{\tau,i}
      &\quad\text{は行列 }M(\tau,i,Y)\text{ により定義},\\[2pt]
    \left(\int_{\varphi}^{\Sigma}Y\right)^{\tau}
      &=\bigoplus_{\varphi(\sigma)=\tau}Y^{\sigma},\\
    d_{\int_{\varphi}^{\Sigma}Y}^{\tau,i}
      &\quad\text{は行列 }M(\tau,i,Y)\text{ により定義}.
  \end{aligned}
  \tag{2.2.3.3}
\end{equation}

\VerdierSemanticAnchor{ch1-2-2-3-pi-extension}
関手 $\int_{\varphi}^{\Pi}$ を関手 $\int_\varphi$ の
$\Pi$-\emph{拡張}と呼ぶ。

\VerdierSemanticAnchor{ch1-2-2-3-sigma-extension}
関手 $\int_{\varphi}^{\Sigma}$ を関手 $\int_\varphi$ の
$\Sigma$-\emph{拡張}と呼ぶ。これら二つの拡張を以下で用いる。

\VerdierSemanticAnchor{ch1-2-2-3-natural-map}
さらに、直和から積への関手的な射に対応する関手の自然な射
$\int_{\varphi}^{\Sigma}\to\int_{\varphi}^{\Pi}$ がある。

\VerdierSemanticAnchor{ch1-2-2-4}
\noindent\textbf{2.2.4.}\quad
$[n]\xrightarrow{\varphi}[m]\xrightarrow{\psi}[r]$ を二つの写像とし、
話を具体的にするため $C$ において可算積が表現可能であると仮定する。
このとき次の関係がある：

\VerdierSemanticAnchor{ch1-2-2-4-eq1}
\begin{equation}
  \int_{\psi}^{\Pi}\int_{\varphi}^{\Pi}
  =\int_{\psi\varphi}^{\Pi},
  \tag{2.2.4.1}
\end{equation}

\VerdierSemanticAnchor{ch1-2-2-4-eq2}
\begin{equation}
  \int_{\operatorname{id}}^{\Pi}=\operatorname{id}.
  \tag{2.2.4.2}
\end{equation}

\VerdierSemanticAnchor{ch1-2-2-4-editorial-caveat}
ここで編者は、いつもの厄介な手品を帽子から取り出すという苦しい義務を負う。
公式 (2.2.4.1) の記号 $=$ が有効なのは、適切な結合律の性質をもつ積関手が
ある場合だけだが、実際にはそのようなことは決してない。

\VerdierSemanticAnchor{ch1-2-2-4-canonical-isomorphism}
一般には、記号「$=$」を (1.2.1.3) 型の通常の整合条件を満たす標準同型で
置き換えなければならない。

\VerdierSemanticAnchor{ch1-2-2-4-sigma-analogue}
関手 $\int_{\varphi}^{\Sigma}$ についても (2.2.4.1) と同様の関係が成り立つ。

\VerdierSemanticAnchor{ch1-2-2-4-summable-restriction}
可算積と直和が必ずしも表現可能でない場合でも、$\psi\varphi$-総和可能な
複体に制限すれば公式 (2.2.4.1) は成り立つ。

\VerdierSemanticAnchor{ch1-2-2-4-simple-complex}
$\psi$ が唯一の射 $[n]\to[1]$ であるとき、記号 $\int_n$ を用いる。
任意の総和可能な $n$-重複複体 $Y$ に対し、複体 $\int_nY$ を
\emph{$Y$ に付随する単純複体}と呼ぶ。
